\SourcePage{375}
\SGAchapter{VI}{射影丛的 $K^\bullet$：计算及其推论}{Le $K^\bullet$ d'un fibré projectif : calculs et conséquences}
\begin{center}Pierre BERTHELOT\end{center}
\SGAsection{1}{射影丛的 $K^\bullet$ 的计算：有限型局部自由层的情形}
本节中，记号 $K^\bullet(X)$ 表示\Term{scheme}{概形} $X$ 上\Term{finite-type}{有限型}\Term{locally-free-sheaf}{局部自由层}的\Term{grothendieck-group}{Grothendieck 群}。本节的目的是证明下述\Term{theorem}{定理}。
\begin{SGAresult}{\Term{theorem}{定理}}{1.1}
设 $S$ 是\Term{quasi-compact}{拟紧}\Term{scheme}{概形}，$E$ 是\Term{rank}{秩}为 $r+1$ 的\Term{locally-free}{局部自由} $\mathcal O_S$-\Term{module}{模}，$X=\mathbb P(E)$ 是相应的\Term{projective-bundle}{射影丛}。\Term{structure-morphism}{结构态射} $f:X\to S$ 定义一个\Term{ring-homomorphism}{环同态}
\[
K^\bullet(S)\xrightarrow{f^*}K^\bullet(X),
\]
从而赋予 $K^\bullet(X)$ 一个 $K^\bullet(S)$-\Term{algebra}{代数}结构。设 $\xi\in K^\bullet(X)$ 为\Term{sheaf}{层} $\mathcal O_X(1)$ 的类。则 $K^\bullet(X)$ 是一个以 $1,\ldots,\xi^r$ 为\Term{basis}{基}的\Term{free-module}{自由} $K^\bullet(S)$-\Term{module}{模}。
\end{SGAresult}
\textbf{证明的安排。}

1）$\xi$ 的幂作为 $K^\bullet(S)$-\Term{module}{模}的\Term{generator}{生成元}生成 $K^\bullet(X)$。为此，我们将证明以下各点。

 a）若 $F$ 是\Term{finite-type}{有限型}\Term{locally-free}{局部自由} $\mathcal O_X$-\Term{module}{模}，则存在一个\Term{finitely-presented}{有限呈示}的\Term{graded-module}{分次} $\operatorname{S}(E)$-\Term{module}{模} $M$，它在 $\mathcal O_S$ 上\Term{flat}{平坦}，并满足 $F=\operatorname{Proj}_0(M)$。这里 $\operatorname{S}(E)$ 是 $E$ 的\Term{symmetric-algebra}{对称代数}；记号 $\operatorname{Proj}_0$ 见 EGA II 8.12.1。

 b）$M$ 有一个\Term{finite}{有限}的\Term{graded-resolution}{分次消解}
\[
0\longrightarrow L_{r+1}\longrightarrow\cdots\longrightarrow L_0\longrightarrow M\longrightarrow0,
\]
其中各 $L_i$ 都是\Term{finite-type}{有限型}\Term{locally-free}{局部自由}的\Term{graded-module}{分次} $\operatorname{S}(E)$-\Term{module}{模}。

 c）各 $L_i$ 有一个\Term{finite}{有限}\Term{filtration}{滤过}，其各\Term{quotient-module}{商}形如 $N_j\otimes_{\mathcal O_S}\operatorname{S}(E)(j)$，其中 $N_j$ 在 $S$ 上\Term{locally-free}{局部自由}且为\Term{finite-type}{有限型}。
\SourcePage{376}
2）$1,\ldots,\xi^{r+1}$ 在 $K^\bullet(S)$ 上\Term{linearly-dependent}{线性相关}。

3）$1,\ldots,\xi^r$ 在 $K^\bullet(S)$ 上\Term{linearly-independent}{线性无关}。

\textbf{1.2。} 为方便起见，先回顾\Term{graded-module}{分次模}的若干经典定义和性质。

设 $S$ 为\Term{scheme}{概形}，$B$ 为\Term{graded-algebra}{分次} $\mathcal O_S$-\Term{algebra}{代数}，$M$ 为\Term{graded-module}{分次} $B$-\Term{module}{模}。对任意\Term{graded-module}{分次} $B$-\Term{module}{模} $N$，照常记 $N(n)$ 为由
\[
(N(n))_k=N_{n+k}\qquad\text{（对任意 $k$）}
\]
定义的\Term{graded-module}{分次} $B$-\Term{module}{模}。称 $M$ 为\Term{locally-free}{局部自由}（相应地，\Term{finite-type}{有限型}\Term{locally-free}{局部自由}）\Term{graded-module}{分次} $B$-\Term{module}{模}，如果 $S$ 的每一点都有一个\Term{open-neighbourhood}{开邻域} $U$，使 $M|U$ 作为\Term{graded-module}{分次} $B|U$-\Term{module}{模}\Term{isomorphism}{同构}于形如 $B(n)|U$ 的\Term{module}{模}的\Term{direct-sum}{直和}（相应地，\Term{finite}{有限}\Term{direct-sum}{直和}）。称 $M$ 为\Term{finitely-presented}{有限呈示}的\Term{graded-module}{分次} $B$-\Term{module}{模}，如果 $S$ 的每一点都有一个\Term{open-neighbourhood}{开邻域} $U$，其上存在\Term{exact-sequence}{正合列} $P\to Q\to M|U\to0$，其中 $P,Q$ 为\Term{finite-type}{有限型}\Term{free-module}{自由}\Term{graded-module}{分次} $B|U$-\Term{module}{模}，各\Term{homomorphism}{同态}均与\Term{grading}{分次}相容。

\textbf{1.2.1。} 对每个整数 $n$，可定义一个\Term{canonical}{典范}的\Term{graded-module}{分次} $B$-\Term{module}{模}\Term{homomorphism}{同态}
\[
M_n\otimes_{\mathcal O_S}B(-n)\longrightarrow M,\tag{1.2.1}
\]
它将 $x\otimes a$ 映到 $M_k$ 中的元素 $a\cdot x$，其中 $x$ 为 $M_n$ 的\Term{section}{截面}，$a$ 为 $(B(-n))_k=B_{k-n}$ 的\Term{section}{截面}。

\textbf{1.2.2。} 若\Term{graded-module}{分次} $B$-\Term{module}{模} $M$ 作为\Term{graded-module}{分次} $B$-\Term{module}{模}\Term{isomorphism}{同构}于某个 $(B(-n))^k$，便称 $M$ 为\Term{finite-type}{有限型}\Term{free-module}{自由}且\Term{homogeneous-module}{齐次}，\Term{degree}{次数}为 $n$。
\SourcePage{377}
设 $M$ 是\Term{finite-type}{有限型}\Term{free-module}{自由}\Term{graded-module}{分次} $B$-\Term{module}{模}。由 1.2，可取一个\Term{graded-isomorphism}{分次同构}
\[
M\simeq\bigoplus_{n=n_0}^{n_1}(B(-n))^{k_n}.\tag{1.2.2}
\]
由此得到 $M$ \Term{resolution}{分解}为\Term{finite-type}{有限型}\Term{free-module}{自由}\Term{homogeneous-module}{齐次}\Term{graded-module}{分次} $B$-\Term{module}{模}的\Term{direct-sum}{直和}：
\[
M\simeq\bigoplus_{n=n_0}^{n_1}M^{(n)},
\]
其中 $M^{(n)}$ 为\Term{degree}{次数} $n$、\Term{rank}{秩} $k_n$ 的\Term{free-module}{自由}\Term{homogeneous-module}{齐次模}。
\begin{SGAresult}{命题}{1.2.3}
\Term{resolution}{分解}（1.2.2）中的整数 $k_n$ 与所选\Term{resolution}{分解}无关。
\end{SGAresult}
\begin{proof}
设给定两个\Term{resolution}{分解}，从而有一个\Term{graded-isomorphism}{分次同构}
\[
\bigoplus_{n=n_0}^{n_1}(B(-n))^{k_n}\simeq
\bigoplus_{n=n'_0}^{n'_1}(B(-n))^{k'_n}.
\]
对左边\Term{resolution}{分解}中非零\Term{free-module}{自由}\Term{homogeneous-module}{齐次}\Term{direct-summand}{直和因子}的个数作归纳。若此数为零，则 $M=0$，命题显然成立。假设已对 $n_1-n_0=m$ 证明，现证明 $n_1-n_0=m+1$ 的情形。给定的\Term{isomorphism}{同构}与\Term{grading}{分次}相容，故在两边的\Term{homogeneous-component}{齐次分量}上诱导\Term{isomorphism}{同构}。特别地，考察\Term{degree}{次数}小于 $n_0$ 的\Term{homogeneous-component}{齐次分量}可知，当 $n<n_0$ 时 $k'_n=0$；考察\Term{degree}{次数} $n_0$ 的\Term{homogeneous-component}{齐次分量}可知，两者的\Term{rank}{秩}相等，故 $k_{n_0}=k'_{n_0}$。

因此，给定的\Term{isomorphism}{同构}在两个\Term{resolution}{分解}的最低次\Term{free-module}{自由}\Term{homogeneous-module}{齐次}\Term{direct-summand}{直和因子}之间诱导\Term{isomorphism}{同构}，从而可对这些\Term{direct-summand}{直和因子}取\Term{quotient-module}{商}。再对所得\Term{quotient-module}{商模}应用归纳假设，即得结论。
\end{proof}
\SourcePage{378}
由 1.2.3 特别可知，对于每个\Term{finite-type}{有限型}\Term{free-module}{自由}\Term{graded-module}{分次模} $M$，它的\Term{homogeneous-basis}{齐次基}中具有给定\Term{degree}{次数}的元\Term{prime-number}{素数}目与所选的\Term{basis}{基}无关。特别地，\Term{resolution}{分解}（1.2.2）中由条件 $k_{n_0}\ne0$、$k_{n_1}\ne0$ 确定的整数 $n_0,n_1$ 与\Term{resolution}{分解}无关。
\begin{SGAresult}{命题}{1.3}
设 $S$ 是\Term{quasi-compact}{拟紧}\Term{scheme}{概形}，$E$ 是\Term{finite-type}{有限型}\Term{locally-free}{局部自由} $\mathcal O_S$-\Term{module}{模}，$X$ 是相应的\Term{projective-bundle}{射影丛}。设 $F$ 是\Term{finitely-presented}{有限呈示}的 $\mathcal O_X$-\Term{module}{模}，在 $S$ 上\Term{flat}{平坦}。存在整数 $n_0$，使得对 $n\geq n_0$ 有：
\begin{enumerate}
\item[i)] $R^if_*(F(n))=0$，$i\geq1$；
\item[ii)] $f_*(F(n))$ 在 $S$ 上\Term{locally-free}{局部自由}且为\Term{finite-type}{有限型}；
\item[iii)] $M=\bigoplus_{n\geq n_0}f_*(F(n))$ 是\Term{finitely-presented}{有限呈示}的\Term{graded-module}{分次} $\operatorname{S}(E)$-\Term{module}{模}（1.2）。
\end{enumerate}
\end{SGAresult}
\begin{proof}
所求性质在 $S$ 上是\Term{local-scheme}{局部}的；因 $S$ \Term{quasi-compact}{拟紧}，取所求各整数的最大值，便可归结到 $S$ \Term{affine}{仿射}的情形。

先设 $S$ 为\Term{noetherian}{诺特}的。性质 i）由 Serre \Term{theorem}{定理}得到。选 $n_0$ 使性质 i）成立，则性质 ii）由 EGA III 7.9.10 得到。

由 EGA II 3.4.5，存在一个\Term{finite-type}{有限型}\Term{quasi-coherent}{拟凝聚}\Term{graded-module}{分次} $\operatorname{S}(E)$-\Term{module}{模} $M'$，使得 $F=\operatorname{Proj}_0(M')$。另一方面，由 EGA III 2.3.1，
\[
\Gamma_*(F)=\Gamma_*(\operatorname{Proj}_0M')
=\bigoplus_{n\in\mathbb Z}f_*(F(n))
\]
与 $M'$ 是 TN-\Term{isomorphism}{同构}的。因此存在整数 $n_1$，使
\[
\bigoplus_{n\geq n_1}f_*(F(n))\simeq\bigoplus_{n\geq n_1}M'_n.
\]
\SourcePage{379}
故可使 $\bigoplus_{n\geq n_1}f_*(F(n))$ 为\Term{finite-type}{有限型} $\operatorname{S}(E)$-\Term{module}{模}；由\Term{noetherian}{诺特}假设，它也为\Term{finitely-presented}{有限呈示}。事实上，这一结论对每个整数 $n_1$ 均成立，因为 Serre \Term{theorem}{定理}保证各 $f_*(F(n))$ 都\Term{coherent}{凝聚}。这样，命题在 $S$ 为\Term{noetherian}{诺特}时已得证。

现设有\Term{cartesian-square}{笛卡尔方块}
\[
\begin{tikzcd}
X\arrow[r,"g"]\arrow[d,"f"']&X_0\arrow[d,"f_0"]\\
S\arrow[r,"g_0"']&S_0,
\end{tikzcd}
\]
其中 $S_0$ 为\Term{noetherian}{诺特}的，$X_0=\mathbb P(E_0)$；设 $F_0$ 是\Term{finitely-presented}{有限呈示}的 $\mathcal O_{X_0}$-\Term{module}{模}，在 $S_0$ 上\Term{flat}{平坦}，并令 $F=g^*(F_0)$。由上文，命题对 $F_0$ 成立。由 EGA III 7.7.5 和 7.8.5，在上述假设下\Term{base-change}{基变换}\Term{theorem}{定理}成立；立即\Term{pushout}{推出}命题对 $F$ 成立。

所以，要完成证明，只须说明\Term{generic}{一般}情形（其中 $S$ \Term{affine}{仿射}）可由\Term{noetherian}{诺特}情形通过\Term{base-change}{基变换}得到。

因 $S$ \Term{affine}{仿射}，它是\Term{noetherian}{诺特}\Term{affine}{仿射}\Term{scheme}{概形}的\Term{filtered-projective-limit}{滤过射影极限}（若 $S=\operatorname{Spec}(A)$，将 $A$ 写成其在 $\mathbb Z$ 上\Term{finite-type}{有限型}的子\Term{algebra}{代数}的\Term{inductive-limit}{归纳极限}即可）。设 $S=\varprojlim S_\lambda$。由 EGA IV 8.5.2 和 8.5.5，存在指标 $\lambda$ 及一个\Term{finite-type}{有限型}\Term{locally-free}{局部自由} $\mathcal O_{S_\lambda}$-\Term{module}{模} $E_\lambda$，使 $E=u_\lambda^*(E_\lambda)$，其中 $u_\lambda$ 表示\Term{morphism}{态射} $S\to S_\lambda$，$u_{\lambda\mu}$ 表示过渡\Term{morphism}{态射} $S_\mu\to S_\lambda$。

对 $\mu\geq\lambda$，令 $E_\mu=u_{\lambda\mu}^*(E_\lambda)$、$X_\mu=\mathbb P(E_\mu)$，则对每个 $\mu$ 有 $X=X_\mu\times_{S_\mu}S$，且 $X=\varprojlim X_\mu$。因此存在指标 $\lambda'\geq\lambda$
\SourcePage{380}
及 $X_{\lambda'}$ 上的\Term{finitely-presented}{有限呈示}\Term{sheaf}{层} $F_{\lambda'}$，使 $F=u_{\lambda'}'^*(F_{\lambda'})$，其中 $u'_{\lambda'}$ 为\Term{morphism}{态射} $X\to X_{\lambda'}$。因各 $X_\mu$ 及 $F_\mu$ 为\Term{finitely-presented}{有限呈示}，由 EGA IV 11.2.6 可取指标 $\lambda''$，使 $F_{\lambda''}$ 在 $S_{\lambda''}$ 上\Term{flat}{平坦}。这样，1.3 的所有假设均满足，证明完成。
\end{proof}
\begin{SGAresult}{推论}{1.4}
在 1.3 的假设下，存在一个\Term{finitely-presented}{有限呈示}的\Term{graded-module}{分次} $\operatorname{S}(E)$-\Term{module}{模} $M$，它在 $\mathcal O_S$ 上\Term{flat}{平坦}，且 $F=\operatorname{Proj}_0(M)$。
\end{SGAresult}
\begin{proof}
取整数 $n_0$ 使 1.3 的性质成立。令 $M=\bigoplus_{n\geq n_0}f_*(F(n))$，则已知 $M$ 为\Term{finitely-presented}{有限呈示}；又因 $n\geq n_0$ 时 $f_*(F(n))$ 在 $S$ 上\Term{locally-free}{局部自由}，故 $M$ 在 $S$ 上\Term{flat}{平坦}。此外，由 EGA II 3.4，$F=\operatorname{Proj}_0(M)$。
\end{proof}
现在进行证明的 b）部分。
\begin{SGAresult}{\Term{lemma}{引理}}{1.5.1}
设 $A$ 为\Term{ring}{环}，$B$ 为非负次的\Term{graded-algebra}{分次} $A$-\Term{algebra}{代数}，且 $B_0=A$；设 $M$ 为\Term{graded-module}{分次} $B$-\Term{module}{模}。假设 $M$ 作为非\Term{graded-module}{分次} $B$-\Term{module}{模}是\Term{finite-type}{有限型}\Term{projective-module}{投射模}。设 $\widetilde B$ 是 $B$ 在 $\operatorname{Spec}(A)$ 上定义的\Term{algebra-sheaf}{代数层}，$\widetilde M$ 为 $M$ 所对应的 $\widetilde B$-\Term{module}{模}。则 $\widetilde M$ 作为\Term{graded-module}{分次} $\widetilde B$-\Term{module}{模}（1.2），在 $\operatorname{Spec}(A)$ 上\Term{local-scheme}{局部}为\Term{finite-type}{有限型}\Term{free-module}{自由模}。
\end{SGAresult}
\begin{proof}
因 $M$ 为\Term{finite-type}{有限型}，它的\Term{degree}{次数}有下界。此外，$M\otimes_BA$ 是\Term{finite-type}{有限型}\Term{projective-module}{投射} $A$-\Term{module}{模}，故只有\Term{finite}{有限}个非零\Term{homogeneous-component}{齐次分量}，且这些\Term{set-component}{分量}本身均为\Term{finite-type}{有限型}\Term{projective-module}{投射} $A$-\Term{module}{模}。在 $\operatorname{Spec}(A)$ 上\Term{localization}{局部化}后，可设 $M\otimes_BA$ 为\Term{finite-type}{有限型}\Term{free-module}{自由} $A$-\Term{module}{模}。最后，因 $M$ \Term{projective-module}{投射}，有 $\operatorname{Tor}_1^B(M,A)=0$。由这些性质可知，$M$ 作为\Term{graded-module}{分次模}是\Term{finite-type}{有限型}\Term{free-module}{自由} $B$-\Term{module}{模}（见 Bourbaki，Alg. 第 II 章 §11，命题 7），故 $\widetilde M$ 是\Term{finite-type}{有限型}\Term{free-module}{自由} $\widetilde B$-\Term{module}{模}，证毕。
\end{proof}
\SourcePage{381}
\begin{SGAresult}{\Term{lemma}{引理}}{1.5.2}
设 $A$ 为\Term{ring}{环}，$B$ 为\Term{finitely-presented}{有限呈示}的 $A$-\Term{algebra}{代数}，$M$ 为\Term{finitely-presented}{有限呈示}的 $B$-\Term{module}{模}。假设 $B,M$ 在 $A$ 上\Term{flat}{平坦}，并且对 $A$ 的每个\Term{residue-field}{剩余域} $k$，$M\otimes_Ak$ 在 $B\otimes_Ak$ 上的\Term{projective-dimension}{投射维数}不超过 $n$（当 $B=A[t_0,\ldots,t_r]$ 时即属此情形）。则 $M$ 有一个由\Term{finite-type}{有限型}\Term{projective-module}{投射} $B$-\Term{module}{模}组成的\Term{length}{长度}为 $n$ 的\Term{resolution}{消解}：
\[
0\longrightarrow L_n\longrightarrow\cdots\longrightarrow L_0\longrightarrow M\longrightarrow0.
\]
\end{SGAresult}
\begin{proof}本\Term{lemma}{引理}立即由 EGA IV 12.3.2 得到。\end{proof}
\begin{SGAresult}{命题}{1.6}
设 $S$ 是\Term{quasi-compact}{拟紧}\Term{scheme}{概形}，$E$ 是\Term{rank}{秩}为 $r+1$ 的\Term{locally-free}{局部自由} $\mathcal O_S$-\Term{module}{模}，$B$ 为其\Term{symmetric-algebra}{对称代数}，赋予通常的\Term{grading}{分次}，$M$ 是\Term{finitely-presented}{有限呈示}的\Term{graded-module}{分次} $B$-\Term{module}{模}，在 $S$ 上\Term{flat}{平坦}。则存在\Term{graded-resolution}{分次消解}
\[
0\longrightarrow L_{r+1}\longrightarrow\cdots\longrightarrow L_0\longrightarrow M\longrightarrow0,
\]
其中各 $L_i$ 均为\Term{graded-module}{分次} $B$-\Term{module}{模}，并且作为\Term{graded-module}{分次} $B$-\Term{module}{模}\Term{locally-free}{局部自由}且为\Term{finite-type}{有限型}。
\end{SGAresult}
\begin{proof}
先证明：每个满足命题条件的\Term{graded-module}{分次} $B$-\Term{module}{模} $M$ 都有一个\Term{surjective}{满}的\Term{graded-homomorphism}{分次同态} $L\to M$，其中 $L$ 是\Term{finite-type}{有限型}\Term{locally-free}{局部自由} $B$-\Term{module}{模}。

因 $M$ 为\Term{finite-type}{有限型}，存在整数 $N$，使 $M$ 由其\Term{degree}{次数}不超过 $N$ 的\Term{set-component}{分量}在 $B$ 上生成。令
\[
L=\bigoplus_{i\leq N}M_i\otimes_{\mathcal O_S}B(-i).
\]
由（1.2.1）定义从 $L$ 到 $M$ 的\Term{homomorphism}{同态}，它显然为\Term{surjective}{满射}。

由 1.5.2，$M$ 在每个\Term{affine-open}{仿射开集}上都有由\Term{finite-type}{有限型}\Term{locally-free}{局部自由} $B$-\Term{module}{模}组成的（非\Term{graded-resolution}{分次}）\Term{resolution}{消解}。因此，若
\SourcePage{382}
$M$ 有一个全局的部分\Term{resolution}{消解}
\[
L_i\longrightarrow L_{i-1}\longrightarrow\cdots\longrightarrow L_0\longrightarrow M\longrightarrow0,
\]
其中各项为\Term{finite-type}{有限型}\Term{locally-free}{局部自由} $B$-\Term{module}{模}，则 $L_i\to L_{i-1}$ 的\Term{kernel}{核} $Z_i$ 仍为\Term{finite-type}{有限型} $B$-\Term{module}{模}（第 I 讲 2.11）。它甚至是\Term{finitely-presented}{有限呈示}的（以 $i+1$ 代替 $i$ 作相同论证即可）。最后，它在 $S$ 上\Term{flat}{平坦}。因此它满足与 $M$ 相同的假设，由开头的论述，就能逐步构造 $M$ 的一个全局\Term{graded-resolution}{分次消解}，各项均为\Term{finite-type}{有限型}\Term{locally-free}{局部自由}\Term{graded-module}{分次} $B$-\Term{module}{模}。当 $i\geq r+1$ 时，$Z_i$ 在每个\Term{affine-open}{仿射开集}上都是\Term{projective-module}{投射} $B$-\Term{module}{模}（1.5.2），故由 1.5.1，它作为\Term{graded-module}{分次} $B$-\Term{module}{模}\Term{locally-free}{局部自由}。于是在所构造的\Term{graded-resolution}{分次消解}中，以 $Z_{r+1}$ 代替 $L_{r+1}$，以 $0$ 代替以后的各项，就得到所述\Term{resolution}{消解}。
\end{proof}
最后证明 c）部分。
\begin{SGAresult}{命题}{1.7}
设 $S$ 是\Term{quasi-compact}{拟紧}\Term{scheme}{概形}，$E$ 为\Term{finite-type}{有限型}\Term{locally-free}{局部自由} $\mathcal O_S$-\Term{module}{模}，$B=\operatorname{S}(E)$，$L$ 是\Term{finite-type}{有限型}\Term{locally-free}{局部自由}\Term{graded-module}{分次} $B$-\Term{module}{模}（1.2）。则 $L$ 有一个由\Term{graded-submodule}{分次子模} $L^{(m)}$ 组成的\Term{descent}{下降}\Term{filtration}{滤过}，使得：
\begin{enumerate}
\item[i)] $m\leq m_0$ 时 $L^{(m)}=L$；
\item[ii)] $m\geq m_1$ 时 $L^{(m)}=0$；
\item[iii)] 对每个 $m$，$\operatorname{gr}^m(L)=L^{(m)}/L^{(m+1)}$ 作为\Term{graded-module}{分次} $B$-\Term{module}{模}\Term{isomorphism}{同构}于形如 $N_m\otimes_{\mathcal O_S}B(m)$ 的\Term{graded-module}{分次} $B$-\Term{module}{模}，其中 $N_m$ 为\Term{finite-type}{有限型}\Term{locally-free}{局部自由} $\mathcal O_S$-\Term{module}{模}。
\end{enumerate}
\end{SGAresult}
\begin{proof}
因 $S$ \Term{quasi-compact}{拟紧}，可取整数 $n_0,n_1$，使 $S$ 每一点的\Term{neighborhood}{邻域}上，$L$ 的\Term{local-scheme}{局部}\Term{generator}{生成元}的\Term{degree}{次数}
\SourcePage{383}
均在 $n_0$ 与 $n_1$ 之间。对 $N=n_1-n_0$ 作归纳。

当 $N=0$，即 $n_1=n_0$ 时，以 $n$ 表示这个共同值。\Term{canonical}{典范}\Term{homomorphism}{同态}（1.2.1）
\[
M_n\otimes_{\mathcal O_S}B(-n)\longrightarrow M
\]
此时是\Term{isomorphism}{同构}。事实上，这是一个\Term{local-scheme}{局部}问题，由（1.2.3）可设 $M\simeq(B(-n))^k$，此时结论显然。

当 $N>0$ 时，再考虑\Term{homomorphism}{同态}（1.2.1）：
\[
M_{n_0}\otimes_{\mathcal O_S}B(-n_0)\longrightarrow M.
\]
\Term{local-scheme}{局部}可知此\Term{homomorphism}{同态}是\Term{injective}{单射}，其\Term{cokernel}{余核} $M'$ 为\Term{finite-type}{有限型}\Term{locally-free}{局部自由}，且\Term{local-scheme}{局部}\Term{generator}{生成元}的\Term{degree}{次数}在 $n_0+1$ 与 $n_1$ 之间。事实上，在 $S$ 的每一点附近，$M_{n_0}$ 或者为零，或者是 $M$ 的最低次非零\Term{homogeneous-component}{齐次分量}。用形如（1.2.2）的\Term{resolution}{分解}将 $M$ \Term{resolution}{分解}为\Term{direct-summand}{直和因子}，并在后一情形应用（1.2.3），可知 $M_{n_0}\otimes_{\mathcal O_S}B(-n_0)$ \Term{isomorphism}{同构}于 $M$ 的\Term{direct-summand}{直和因子} $(B(-n_0))^{k_{n_0}}$，故有上述性质。对 $M'$ 应用归纳假设，立即得到所求结论。
\end{proof}
\textbf{1.8。} 回到\Term{theorem}{定理} 1.1 的证明，现证明 $\xi$ 的幂作为 $K^\bullet(S)$-\Term{module}{模}的\Term{generator}{生成元}生成 $K^\bullet(X)$。

若 $F$ 为\Term{finite-type}{有限型}\Term{locally-free}{局部自由} $\mathcal O_X$-\Term{module}{模}，由 1.4，存在一个\Term{finitely-presented}{有限呈示}的\Term{graded-module}{分次} $B$-\Term{module}{模} $M$，它在 $S$ 上\Term{flat}{平坦}，且 $F=\operatorname{Proj}_0(M)$。由 1.6，$M$ 有一个\Term{graded-resolution}{分次消解}，各项为\Term{finite-type}{有限型}\Term{locally-free}{局部自由}\Term{graded-module}{分次} $B$-\Term{module}{模}。由此得到\Term{exact-sequence}{正合列}
\SourcePage{384}
\[
0\longrightarrow\operatorname{Proj}_0(L_{r+1})\longrightarrow\cdots
\longrightarrow\operatorname{Proj}_0(L_0)\longrightarrow\operatorname{Proj}_0(M)\longrightarrow0,
\]
以及 $K^\bullet(X)$ 中的等式
\[
\operatorname{cl}(F)=\sum_i(-1)^i\operatorname{cl}(\operatorname{Proj}_0(L_i)).
\]
利用 1.7 定义的 $L_i$ 上的\Term{filtration}{滤过}，得
\[
\operatorname{cl}(\operatorname{Proj}_0(L_i))
=\sum_j\operatorname{cl}\bigl(\operatorname{Proj}_0(N_{ij}\otimes_{\mathcal O_S}B(j))\bigr).
\]
而 $\operatorname{Proj}_0(N_{ij}\otimes_{\mathcal O_S}B(j))=f^*(N_{ij})(j)$，故
\[
\begin{aligned}
\operatorname{cl}(\operatorname{Proj}_0(L_i))
&=\sum_j\operatorname{cl}(f^*(N_{ij}))\cdot\xi^j\\
&=\sum_j f^*(\operatorname{cl}(N_{ij}))\cdot\xi^j.
\end{aligned}
\]
最终便将 $\operatorname{cl}(F)$ 表成了 $\xi$ 的幂以 $K^\bullet(S)$ 中元素为系数的\Term{linear-combination}{线性组合}。

\textbf{1.9。} 要证明 $1,\xi,\ldots,\xi^{r+1}$ 在 $K^\bullet(S)$ 上\Term{linearly-dependent}{线性相关}，考虑\Term{canonical}{典范}\Term{surjective}{满}\Term{homomorphism}{同态} $f^*(E)\to\mathcal O_X(1)$，其\Term{kernel}{核} $F$ \Term{locally-free}{局部自由}，\Term{rank}{秩}为 $r$。因而可以写成
\[
\lambda^{r+1}(f^*(E)-\xi)=0,\qquad E=\operatorname{cl}E,
\]
这里 $K^\bullet(X)$ 上的运算 $\lambda^i$ 已在第 V 讲 2.2 中定义。于是
\[
\sum_{i=0}^{r+1}\lambda^i(f^*(E))\lambda^{r+1-i}(-\xi)=0.
\]
因 $\mathcal O_X(1)$ 为\Term{invertible-sheaf}{可逆层}，有 $\lambda_t(\xi)=1+\xi t$，从而 $\lambda_t(-\xi)=1/(1+\xi t)$，最后得到关系
\[
\sum_{i=0}^{r+1}(-1)^{r+1-i}\xi^{r+1-i}\lambda^i(E)=0.\tag{1.9.1}
\]
最后证明 $1,\xi,\ldots,\xi^r$ 在 $K^\bullet(S)$ 上\Term{linearly-independent}{线性无关}。\footnote{这一证明出自 J. P. Jouanolou。}
\SourcePage{385}
\begin{SGAresult}{\Term{lemma}{引理}}{1.10}
沿用 1.1 的记号与假设，设 $F,G$ 是 $X$ 上的两个\Term{finite-type}{有限型}\Term{locally-free-sheaf}{局部自由层}，满足
\[
\operatorname{cl}_{K^\bullet(X)}(F)=\operatorname{cl}_{K^\bullet(X)}(G).
\]
则存在整数 $n_0$，使 $n\geq n_0$ 时：
\begin{enumerate}
\item[i)] $f_*(F(n))$、$f_*(G(n))$ 在 $S$ 上\Term{locally-free}{局部自由}且为\Term{finite-type}{有限型}；
\item[ii)] $\operatorname{cl}_{K^\bullet(S)}(f_*(F(n)))=\operatorname{cl}_{K^\bullet(S)}(f_*(G(n)))$。
\end{enumerate}
\end{SGAresult}
\begin{proof}
设 $A$ 为由 $X$ 上的\Term{finite-type}{有限型}\Term{locally-free-sheaf}{局部自由层}的\Term{isomorphism-class}{同构类}生成的\Term{free-abelian-group}{自由阿贝尔群}。由 $K^\bullet$ 的构造，$F,G$ 在 $K^\bullet(X)$ 中同类，意味着存在\Term{finite-type}{有限型}\Term{locally-free}{局部自由} $\mathcal O_X$-\Term{module}{模}的\Term{exact-sequence}{正合列}
\[
0\longrightarrow L_i\longrightarrow M_i\longrightarrow N_i\longrightarrow0,
\qquad i=1,\ldots,k,
\]
及整数 $n_i$，使 $A$ 中有等式
\[
(F)-(G)=\sum_i n_i\bigl((L_i)-(M_i)+(N_i)\bigr).
\]
对每个 $n$，与 $\mathcal O_X(n)$ 作\Term{tensor-product}{张量积}定义 $A$ 的一个\Term{automorphism}{自同构}，故
\[
(F(n))-(G(n))=\sum_i\bigl((L_i(n))-(M_i(n))+(N_i(n))\bigr),
\]
而各列 $0\to L_i(n)\to M_i(n)\to N_i(n)\to0$ 仍\Term{exact-sequence}{正合}。

设 $B$ 为由 $S$ 上\Term{quasi-coherent}{拟凝聚}\Term{sheaf}{层}的\Term{isomorphism-class}{同构类}生成的\Term{free-abelian-group}{自由阿贝尔群}，$A'\subset B$ 为由 $S$ 上\Term{finite-type}{有限型}\Term{locally-free-sheaf}{局部自由层}的类生成的\Term{subgroup}{子群}。由\Term{direct-image}{正像}定义一个\Term{homomorphism}{同态} $A\to B$，于是还有
\SourcePage{386}
\[
\begin{aligned}
(f_*(F(n)))-(f_*(G(n)))
=\sum_i n_i\bigl(&(f_*(L_i(n)))-(f_*(M_i(n)))\\
&+(f_*(N_i(n)))\bigr).
\end{aligned}
\]
由 1.2，当 $n$ 充分大时，下述条件成立：
\[
\forall k\geq1,\ \forall i,\qquad R^kf_*(L_i(n))=0;
\]
并且对每个 $i$，$f_*(L_i(n))$、$f_*(M_i(n))$、$f_*(N_i(n))$、$f_*(F(n))$、$f_*(G(n))$ 均\Term{locally-free}{局部自由}且为\Term{finite-type}{有限型}。

因此，当 $n$ 充分大时，这些元素都在 $A'$ 中，满足上面写出的关系，而且各列
\[
0\longrightarrow f_*(L_i(n))\longrightarrow f_*(M_i(n))
\longrightarrow f_*(N_i(n))\longrightarrow0
\]
均\Term{exact-sequence}{正合}。这就说明，$f_*(F(n))$ 与 $f_*(G(n))$ 在 $K^\bullet(S)$ 中同类。
\end{proof}
\textbf{1.11。} 假设 $1,\xi,\ldots,\xi^r$ 之间有一个在 $K^\bullet(S)$ 上的\Term{linear}{线性}关系：
\[
f^*(x_0)+f^*(x_1)\xi+\cdots+f^*(x_r)\xi^r=0.\tag{1.11.1}
\]
每个 $x_j$ 均可写成两个\Term{sheaf}{层}的类之差，即 $x_j=\operatorname{cl}(E_j)-\operatorname{cl}(E'_j)$，其中 $E_j,E'_j$ 在 $S$ 上\Term{locally-free}{局部自由}且为\Term{finite-type}{有限型}。令
\[
F=\bigoplus_{j=0}^r f^*(E_j)(j),\qquad
F'=\bigoplus_{j=0}^r f^*(E'_j)(j).
\]
等式（1.11.1）可写成
\[
\operatorname{cl}_{K^\bullet(X)}(F)=\operatorname{cl}_{K^\bullet(X)}(F').
\]
由\Term{lemma}{引理} 1.10，对充分大的 $n$ 有
\[
\operatorname{cl}_{K^\bullet(S)}(f_*(F(n)))=
\operatorname{cl}_{K^\bullet(S)}(f_*(F'(n))).
\]
\SourcePage{387}
而
\[
\begin{aligned}
f_*(F(n))
&=\bigoplus_j f_*\bigl(f^*(E_j)\otimes_{\mathcal O_X}\mathcal O_X(n+j)\bigr)\\
&=\bigoplus_j E_j\otimes_{\mathcal O_S}f_*(\mathcal O_X(n+j)).
\end{aligned}
\]
此外，$f_*(\mathcal O_X(n+j))=\operatorname{S}(E)^{n+j}$；用 $\sigma_p$ 表示 $\operatorname{S}(E)^p$ 在 $K^\bullet(S)$ 中的类。故关系（1.11.1）最终给出
\[
\forall n\geq n_0,\qquad\sum_{j=0}^r x_j\sigma_{n+j}=0.\tag{1.11.2}
\]
现考虑\Term{complex}{复形}
\[
0\longrightarrow\bigwedge^{r+1}(E)\otimes_{\mathcal O_S}\operatorname{S}(E)
\longrightarrow\bigwedge^r(E)\otimes_{\mathcal O_S}\operatorname{S}(E)
\longrightarrow\cdots\longrightarrow E\otimes_{\mathcal O_S}\operatorname{S}(E)
\longrightarrow\operatorname{S}(E)\longrightarrow0,
\]
其\Term{arrow}{箭头}定义为
\[
\begin{aligned}
{}&(x_1\wedge\cdots\wedge x_n)\otimes(y_1\odot\cdots\odot y_p)\\
&\quad\longmapsto\sum_i(-1)^{i-1}
(x_1\wedge\cdots\wedge\widehat{x_i}\wedge\cdots\wedge x_n)
\otimes(x_i\odot y_1\odot\cdots\odot y_p).
\end{aligned}
\]
若将 $\bigwedge^n(E)$ 看作集中在\Term{degree}{次数} $n$ 的\Term{graded-module}{分次} $\mathcal O_S$-\Term{module}{模}，这便是\Term{graded-module}{分次} $\operatorname{S}(E)$-\Term{module}{模}的\Term{complex}{复形}，且各\Term{homomorphism}{同态}本身也是\Term{graded-homomorphism}{分次}的。

这一\Term{complex}{复形}是 $\mathcal O_S$ 的一个\Term{resolution}{消解}：在 $E$ \Term{free-module}{自由}的\Term{open-subset}{开集}上，$\operatorname{S}(E)$ 可认同为\Term{polynomial-algebra}{多项式代数} $\mathcal O_S[T_0,\ldots,T_r]$，该\Term{complex}{复形}就是由 $\operatorname{S}(E)$-\Term{regular-sequence}{正则序列} $T_0,\ldots,T_r$ 定义的\Term{koszul-complex}{Koszul 复形}；已知后者是 $\operatorname{S}(E)/(T_0,\ldots,T_r)=\mathcal O_S$ 的\Term{resolution}{消解}。

将此\Term{complex}{复形}的每个\Term{degree}{次数} $\geq1$ 的\Term{homogeneous-component}{齐次分量}为\Term{acyclic-complex}{无上同调复形}这一事实写出，即在 $K^\bullet(S)$ 中得到关系
\[
\sum_{i=0}^{r+1}(-1)^i\lambda^i(E)\sigma_{p-i}=0
\qquad\text{对每个 $p\geq1$},\tag{1.11.3}
\]
其中约定 $k<0$ 时 $\sigma_k=0$。

假设关系（1.11.2）对 $n\geq n_0$ 成立，下面证明它对 $n_0-1$ 仍成立，从而可递降归纳至 $n_0=-r$（稍后将看到，这足以\Term{pushout}{推出}各 $x_j$ 均为零）。
\SourcePage{388}
由（1.11.3）可写出
\[
\sum_{j=0}^r x_j\left(\sum_{i=0}^{r+1}(-1)^i\lambda^i(E)\sigma_{p(j)-i}\right)=0,
\]
其中 $p(j)\geq1$。取 $p(j)=n_0+r+j$；只要 $n_0>-r$，这就确实不小于 $1$。于是
\[
\begin{aligned}
\sum_{j=0}^r x_j\left(\sum_{i=0}^{r+1}(-1)^i\lambda^i(E)\sigma_{n_0+r+j-i}\right)
=\sum_{i=0}^{r+1}(-1)^i\lambda^i(E)
\left(\sum_{j=0}^r x_j\sigma_{n_0+r+j-i}\right).
\end{aligned}
\]
当 $i\ne r+1$ 时，$n_0+r-i\geq n_0$，故由（1.11.2），
\[
\sum_{j=0}^r x_j\sigma_{n_0+r+j-i}=0.
\]
因此，关系（1.11.4）化为
\[
\lambda^{r+1}(E)\left(\sum_{j=0}^r x_j\sigma_{n_0-1+j}\right)=0.
\]
而 $E$ 的\Term{rank}{秩}为 $r+1$，故 $\lambda^{r+1}(E)$ 是一个\Term{invertible-sheaf}{可逆层}的类，从而在 $K^\bullet(S)$ 中\Term{invertible}{可逆}。因此关系（1.11.4）化为 $n_0-1$ 时的关系（1.11.2）。

所以（1.11.2）对 $n\geq-r$ 成立；取 $n=-r$，得
\[
x_r\sigma_0=0.
\]
因 $\sigma_0=1$，得 $x_r=0$。利用 $n$ 从 $-r$ 到 $0$ 的关系（1.11.2），对 $j$ 作递降归纳，便得到各 $x_j$ 全为零。这证明了 $1,\ldots,\xi^r$ 在 $K^\bullet(S)$ 上\Term{linearly-independent}{线性无关}，从而完成\Term{theorem}{定理} 1.1 的证明。

\textbf{注记 1.12。} 对每个\Term{finite-type}{有限型}\Term{locally-free}{局部自由} $\mathcal O_S$-\Term{module}{模} $E$，设 $\sigma_i(E)$ 为 $E$ 的第 $i$ 个\Term{symmetric-power}{对称幂}在 $K^\bullet(S)$ 中的类；令
\[
\sigma_t(E)=1+\sum_{i\geq1}\sigma_i(E)t^i.
\]
关系（1.11.3）还可写成
\SourcePage{389}
\[
\lambda_t(E)\cdot\sigma_{-t}(E)=1.\tag{1.12.1}
\]
这一关系本身表明映射 $\sigma_t$ 的\Term{additivity}{可加性}，因此它可延拓为从 $K^\bullet(S)$ 到 $1+K^\bullet(S)[[t]]^+$ 的\Term{homomorphism}{同态}。

\textbf{注记 1.13。} $\xi$ 的任意一串幂 $\xi^k,\ldots,\xi^{k+r+1}$ 均可用作 $K^\bullet(X)$ 在 $K^\bullet(S)$ 上的\Term{basis}{基}，因为 $\xi$ \Term{invertible}{可逆}，这样一串元素便在 $K^\bullet(S)$ 上\Term{linearly-independent}{线性无关}；另一方面，由等式（1.9.1）及 $\xi$ 的\Term{invertible}{可逆}性，$\xi$ 的其他幂都可显式表示成这些幂以 $K^\bullet(S)$ 中元素为系数的\Term{linear-combination}{线性组合}。

令 $x=1-\xi^{-1}$（$x$ \Term{solvable}{可解}释为“一个\Term{hyperplane-section}{超平面截面}的类”）。则 $1,x,\ldots,x^r$ 也可用作 $K^\bullet(X)$ 在 $K^\bullet(S)$ 上的\Term{basis}{基}。要得到 $1,\ldots,x^{r+1}$ 的相关关系，考虑\Term{exact-sequence}{正合列}
\[
0\longrightarrow F\longrightarrow f^*(E)\longrightarrow\mathcal O_X(1)\longrightarrow0,
\]
以及取\Term{dual}{对偶}所得的\Term{exact-sequence}{正合列}
\[
0\longrightarrow\mathcal O_X(-1)\longrightarrow f^*(E^\vee)\longrightarrow F^\vee\longrightarrow0.
\]
此时注意到 $\gamma^{r+1}(F^\vee-r)=0$，由此\Term{pushout}{推出}
\[
\gamma^{r+1}(F^\vee-r)=\gamma^{r+1}(E^\vee-r-1+x)=0,
\]
亦即
\[
\sum_{k=0}^{r+1}\gamma^{r+1-k}(E^\vee-r-1)\cdot x^k=0.\tag{1.13.1}
\]
引入 $x'=\xi-1$ 可得到类似结果。

最后注意，若 $E=(\mathcal O_S)^{r+1}$，从而 $X=\mathbb P_S^r$，则 $E=r+1$，等式（1.13.1）化为
\SourcePage{390}
\[
x^{r+1}=0.\tag{1.13.2}
\]
\textbf{注记 1.14。} 使用“\Term{relative-scheme}{相对概形}”的概念（参见 [3]），还可在假设 $S$ 为\Term{locally-ringed-topos}{局部环化拓扑斯}时赋予 1.1 意义。若假设此\Term{topos}{拓扑斯}的\Term{final-object}{终对象}\Term{quasi-compact}{拟紧}，则\Term{scheme}{概形}情形中给出的证明可直接移用，\Term{theorem}{定理}因而仍然成立。
\SGAsection{2}{射影丛的 $K^\bullet$ 的计算：完美复形的情形}
本节在本讨论班后续内容中不再使用。

本节中，记号 $K^\bullet(X)$ 表示 $X$ 上\Term{perfect-complex}{完美复形}的\Term{grothendieck-group}{Grothendieck 群}（第 IV 讲 2.2）。我们要证明以下与 1.1 类似的\Term{theorem}{定理}。
\begin{SGAresult}{\Term{theorem}{定理}}{2.1}
设 $S$ 是\Term{quasi-compact}{拟紧}\Term{scheme}{概形}，$E$ 是\Term{rank}{秩}为 $r+1$ 的\Term{locally-free}{局部自由} $\mathcal O_S$-\Term{module}{模}，$X=\mathbb P(E)$ 为相应的\Term{projective-bundle}{射影丛}；设 $\xi$ 是\Term{degree}{次数}零处为 $\mathcal O_X(1)$、其他\Term{degree}{次数}处为零的\Term{complex}{复形}在 $K^\bullet(X)$ 中的类。则 $K^\bullet(X)$ 配备其自然的 $K^\bullet(S)$-\Term{algebra}{代数}结构后，是以 $1,\xi,\ldots,\xi^r$ 为\Term{basis}{基}的\Term{free-module}{自由} $K^\bullet(S)$-\Term{module}{模}。
\end{SGAresult}
\textbf{2.2。} 固定记号：$B$ 表示 $E$ 的\Term{symmetric-algebra}{对称代数}；以 $\mathscr A$ 表示 $\mathcal O_X$-\Term{module}{模}的\Term{category}{范畴}，$\mathscr C$ 表示\Term{graded-module}{分次} $B$-\Term{module}{模}的\Term{category}{范畴}，其\Term{morphism}{态射}为\Term{graded-morphism}{分次态射}；$\mathscr C_0$ 表示 $\mathscr C$ 的\Term{full-subcategory}{全子范畴}，\Term{object}{对象}为作为\Term{graded-module}{分次} $B$-\Term{module}{模}是\Term{finite-type}{有限型}\Term{locally-free}{局部自由}的 $B$-\Term{module}{模}；$\mathscr C'$ 为 $\mathcal O_S$-\Term{module}{模}的\Term{category}{范畴}。\Term{category}{范畴} $\mathscr A,\mathscr C,\mathscr C'$ 均为\Term{abelian-category}{阿贝尔范畴}。

对每个 $n$ 及 $D^+(\mathscr A)$ 中每个\Term{complex}{复形} $L^\bullet$，以 $L^\bullet(n)$ 表示\Term{complex}{复形} $L^\bullet\otimes_{\mathcal O_X}\mathcal O_X(n)$。
\SourcePage{391}
\textbf{2.3。} 我们需要将\Term{functor}{函子} $\operatorname{Proj}_0$ 应用于任意\Term{graded-module}{分次} $B$-\Term{module}{模}，而不仅像 EGA II 3.2 的构造那样仅用于\Term{quasi-coherent}{拟凝聚}\Term{graded-module}{分次} $B$-\Term{module}{模}。因此，下面简要说明\Term{generic}{一般}情形中的构造。

\textbf{2.3.1。} 先设 $S$ 为\Term{scheme}{概形}，$A$ 为\Term{quasi-coherent}{拟凝聚} $\mathcal O_S$-\Term{algebra}{代数}，$Y=\operatorname{Spec}(A)$，$f:Y\to S$ 为\Term{structure-morphism}{结构态射}。由\Term{isomorphism}{同构}
\[
A\xrightarrow{\sim}f_*(\mathcal O_Y)
\]
得到\Term{canonical}{典范}\Term{homomorphism}{同态}
\[
f^{-1}(A)\longrightarrow\mathcal O_Y,\tag{2.3.1.1}
\]
其中 $f^{-1}$ 表示\Term{set}{集合}\Term{sheaf}{层}意义下的\Term{inverse-image}{逆像}；换言之，得到一个\Term{ringed-space}{赋环空间}\Term{morphism}{态射}
\[
\bar f:(Y,\mathcal O_Y)\longrightarrow(S,A).\tag{2.3.1.2}
\]
由此得到从 $A$-\Term{module}{模}\Term{category}{范畴}到 $\mathcal O_Y$-\Term{module}{模}\Term{category}{范畴}的“\Term{inverse-image}{逆像}”\Term{functor}{函子} $\bar f^*$；它是从 $\mathcal O_Y$-\Term{module}{模}\Term{category}{范畴}到 $f_*(\mathcal O_Y)$-\Term{module}{模}\Term{category}{范畴}的“\Term{direct-image}{正像}”\Term{functor}{函子} $\bar f_*$ 的\Term{left-adjoint}{左伴随}。因此，对每个 $A$-\Term{module}{模} $M$，有
\[
\bar f^*(M)=f^{-1}(M)\otimes_{f^{-1}(A)}\mathcal O_Y.\tag{2.3.1.3}
\]
因\Term{homomorphism}{同态}（2.3.1.1）\Term{flat}{平坦}（它在每个\Term{stalk}{茎}上都对应于一个\Term{localization}{局部化}），如此定义的\Term{functor}{函子} $\bar f^*$ 为\Term{exact-functor}{正合函子}（参见 [4]，第 164 页）。

若 $M$ 是\Term{quasi-coherent}{拟凝聚} $A$-\Term{module}{模}，此构造给出 EGA II 1.4.3 中定义的\Term{module}{模} $\widetilde M$。事实上，\Term{isomorphism}{同构}
\SourcePage{392}
\[
M\xrightarrow{\sim}f_*(\widetilde M)=\bar f_*(\widetilde M)
\]
由\Term{adjunction}{伴随}给出一个\Term{homomorphism}{同态}
\[
\bar f^*(M)\longrightarrow\widetilde M,
\]
并且在每个\Term{stalk}{茎}上可立即看出它是\Term{isomorphism}{同构}。

\textbf{2.3.2。} 现设 $B$ 为\Term{quasi-coherent}{拟凝聚}\Term{graded-algebra}{分次} $\mathcal O_S$-\Term{algebra}{代数}，$M$ 为\Term{graded-module}{分次} $B$-\Term{module}{模}，$X=\operatorname{Proj}(B)$。对 $S$ 的每个\Term{affine-open}{仿射开集} $U$，以及 $B$ 在 $U$ 上每个\Term{degree}{次数} $d>0$ 的\Term{homogeneous-section}{齐次截面} $s$，以 $M_{(s),U}$ 表示\Term{fraction-module}{分式模} $M_s$ 的\Term{degree}{次数}零的\Term{homogeneous-component}{齐次分量}；于是
\[
M_{(s),U}=\bigl((M|U)\otimes_{B|U}(B|U)_s\bigr)_0.
\]
若 $V$ 为 $U$ 的\Term{affine-open}{仿射开集}，则 $M_{(s),V}$ 是 $M_{(s),U}$ 在 $V$ 上的\Term{restriction}{限制}。

考虑由下列\Term{affine-open}{仿射开集}构成的 $X$ 的\Term{open-basis}{开集基}：
\[
D_+(U,s)\simeq\operatorname{Spec}((B|U)_{(s)}).
\]
以 $f$ 表示\Term{structure-morphism}{结构态射} $X\to S$，以 $f_{U,s}$ 表示\Term{affine-morphism}{仿射态射}
\[
f_{U,s}:D_+(U,s)\longrightarrow U.
\]
用 2.3.1 的记号，令
\[
\operatorname{Proj}_0(M)_{U,s}=\bar f_{U,s}^*(M_{(s),U}).
\]
我们将证明，当 $U,s$ 变化时，这些 $\operatorname{Proj}_0(M)_{U,s}$ 可以\Term{gluing}{粘合}，从而定义一个记作 $\operatorname{Proj}_0(M)$ 的 $\mathcal O_X$-\Term{module}{模}。固定 $s$ 而令 $U$ 变化时的\Term{gluing}{粘合}显然，因此只须证明（参见 EGA II 2.5.2）：若 $t$ 为 $B$ 在 $U$ 上\Term{degree}{次数} $e>0$ 的\Term{homogeneous-section}{齐次截面}，则 $\operatorname{Proj}_0(M)_{U,s}$ 在 $D_+(U,st)$ 上的\Term{restriction}{限制}\Term{canonical}{典范}\Term{isomorphism}{同构}于 $\operatorname{Proj}_0(M)_{U,st}$。考虑
\SourcePage{393}
\[
\begin{tikzcd}
D_+(st)\arrow[r,hook,"i"]\arrow[dr,"{f_{U,st}}"']&D_+(s)\arrow[d,"{f_{U,s}}"]\\
&U.
\end{tikzcd}
\]
即须证明\Term{module}{模} $i^*\bar f_{U,s}^*(M_{(s),U})$ 与 $\bar f_{U,st}^*(M_{(st),U})$ \Term{canonical}{典范}\Term{isomorphism}{同构}。一方面，有\Term{canonical}{典范}\Term{isomorphism}{同构}
\[
\begin{aligned}
i^*\bar f_{U,s}^*(M_{(s),U})
&\simeq i^*\bigl(f_{U,s}^{-1}(M_{(s),U})
\otimes_{f_{U,s}^{-1}((B|U)_{(s)})}\mathcal O_{D_+(s)}\bigr)\\
&\simeq i^{-1}\bigl(f_{U,s}^{-1}(M_{(s),U})
\otimes_{f_{U,s}^{-1}((B|U)_{(s)})}\mathcal O_{D_+(s)}\bigr)
\otimes_{i^{-1}(\mathcal O_{D_+(s)})}\mathcal O_{D_+(st)}\\
&\simeq f_{U,st}^{-1}(M_{(s),U})
\otimes_{f_{U,st}^{-1}((B|U)_{(s)})}\mathcal O_{D_+(st)}.
\end{aligned}
\]
另一方面，
\[
\bar f_{U,st}^*(M_{(st),U})\simeq
f_{U,st}^{-1}(M_{(st),U})
\otimes_{f_{U,st}^{-1}((B|U)_{(st)})}\mathcal O_{D_+(st)}.
\]
最后，将 EGA II 2.2.2 译为\Term{sheaf}{层}的语言，便得到\Term{canonical}{典范}\Term{isomorphism}{同构}
\[
M_{(st),U}\simeq(M_{(s),U})_{t^d/s^e}
\simeq M_{(s),U}\otimes_{(B|U)_{(s)}}(B|U)_{(st)},
\]
最终得到所求\Term{isomorphism}{同构}。

显然，当 $M$ \Term{quasi-coherent}{拟凝聚}时，如此构造的\Term{module}{模} $\operatorname{Proj}_0(M)$ 正是 EGA II 中定义的那个。

\textbf{2.3.3。} 由 2.3.1、2.3.2 可知，\Term{functor}{函子} $\operatorname{Proj}_0$ 为\Term{exact-functor}{正合}。回到 2.2 的记号，过渡到\Term{derived-category}{导出范畴}，得到仍记作 $\operatorname{Proj}_0$ 的\Term{functor}{函子}
\[
\operatorname{Proj}_0:D^+(\mathscr C)\longrightarrow D^+(\mathscr A).
\]
由以上构造还可看出，
\SourcePage{394}
\Term{graded-module}{分次} $B$-\Term{module}{模}\Term{category}{范畴}上的\Term{functor}{函子} $\operatorname{Proj}_0$ 与\Term{tensor-product}{张量积}相容；所得\Term{derived-category}{导出范畴}上的\Term{functor}{函子}也同样如此。

\textbf{2.4。} 对每个整数 $k$ 及每个 $\mathcal O_X$-\Term{module}{模} $L$，令
\[
F_k(L)=\bigoplus_{n\geq k}f_*(L(n)).\tag{2.4.1}
\]
如此定义的\Term{functor}{函子} $F_k:\mathscr A\to\mathscr C$ 为\Term{left-exact}{左正合}。因 $\mathcal O_X$-\Term{module}{模}\Term{category}{范畴}有足够多的\Term{injective-object}{内射对象}，可定义\Term{derived-functor}{导出函子}
\[
\mathbf RF_k:D^+(\mathscr A)\longrightarrow D^+(\mathscr C).\tag{2.4.2}
\]
若 $L$ 是\Term{finite-type}{有限型}\Term{locally-free}{局部自由} $\mathcal O_X$-\Term{module}{模}，则存在整数 $k_0$，使得对 $k\geq k_0$，$L$ 是 $F_k$-\Term{acyclic-object}{无上同调}的。只须证明，对充分大的 $n$，$L$ 对\Term{functor}{函子} $L\mapsto f_*(L(n))$ 是\Term{acyclic-object}{无上同调}的。因与 $\mathcal O_X(n)$ 作\Term{tensor-product}{张量积}是\Term{exact-functor}{正合函子}，这立即由复合\Term{functor}{函子}的\Term{spectral-sequence}{谱序列}及 1.3 得到。

\textbf{2.5。} 我们要对每个 $\mathcal O_X$-\Term{module}{模} $F$ 构造一个\Term{functorial}{函子性地}依赖于 $F$ 的\Term{homomorphism}{同态}
\[
\beta:\operatorname{Proj}_0(F_k(F))\longrightarrow F.\tag{2.5.1}
\]
为此，定义 $\beta$ 在各 $D_+(U,s)$（2.3.2）上的\Term{restriction}{限制}。令 $M=F_k(F)$，由构造有
\[
\operatorname{Proj}_0(M)|D_+(U,s)=\bar f_{U,s}^*(M_{(s),U}).
\]
要定义\Term{homomorphism}{同态}
\[
\bar f_{U,s}^*(M_{(s),U})\longrightarrow F|D_+(U,s),
\]
\SourcePage{395}
由 $\bar f_{U,s}^*$ 是\Term{direct-image}{正像}\Term{functor}{函子}的\Term{left-adjoint}{左伴随}这一定义，只须定义一个 $f_{U,s*}(\mathcal O_X)$-\Term{linear}{线性}\Term{homomorphism}{同态}
\[
M_{(s),U}\longrightarrow f_{U,s*}(F|D_+(U,s)).
\]
为此，须对 $U$ 的每个\Term{open-subset}{开集} $V$ 给出一个\Term{homomorphism}{同态}
\[
\Gamma(V,M)_{(s)}\longrightarrow\Gamma(D_+(V,s),F),\tag{2.5.2}
\]
使它与对 $U$ 的各\Term{open-subset}{开集}的\Term{restriction}{限制}相容，再取相应的\Term{associated-sheaf}{伴随层}之间的\Term{homomorphism}{同态}。\Term{homomorphism}{同态}（2.5.2）定义如下：设 $z/s^n\in\Gamma(V,M)_{(s)}$，其中 $s$ 的\Term{degree}{次数}为 $d$，
\[
z\in\Gamma(V,M_{nd})=\Gamma(D_+(V,s),F(nd));
\]
将 $z/s^n$ 对应到 $F$ 在 $D_+(V,s)$ 上的\Term{section}{截面} $z\cdot\alpha_d(s^n)^{-1}$，这里 $\alpha_d$ 表示\Term{canonical}{典范}\Term{homomorphism}{同态}
\[
\Gamma(V,B_d)\longrightarrow\Gamma(f^{-1}(V),\mathcal O_X(d))
\]
（参见 EGA II 2.6.3、2.6.4）。所需的相容性容易验证，从而经\Term{gluing}{粘合}得到\Term{homomorphism}{同态}（2.5.1）。

显然，当 $F$ \Term{quasi-coherent}{拟凝聚}时，如此构造的\Term{homomorphism}{同态}正是 EGA II（3.3.4.1）中的那个。

$\mathcal O_X$-\Term{module}{模}\Term{category}{范畴}上的\Term{functor-morphism}{函子同态} $\operatorname{Proj}_0\circ F_k\to\operatorname{Id}$，过渡到\Term{derived-category}{导出范畴}后给出 $D^+(\mathscr A)$ 上的\Term{functor-morphism}{函子同态}
\[
\operatorname{Proj}_0\circ\mathbf RF_k\longrightarrow\operatorname{Id}_{D^+(\mathscr A)}.\tag{2.5.3}
\]
\begin{SGAresult}{命题}{2.6}
设 $S$ 是\Term{quasi-compact}{拟紧}\Term{scheme}{概形}，$E$ 为\Term{finite-type}{有限型}\Term{locally-free}{局部自由} $\mathcal O_S$-\Term{module}{模}，$X=\mathbb P(E)$，$L^\bullet$ 是 $X$ 上的\Term{perfect-complex}{完美复形}。则存在整数 $k_0$，使 $k\geq k_0$ 时以下性质成立：
\SourcePage{396}
\begin{enumerate}
\item[i)] $\mathbf RF_k(L^\bullet)$ 是 $\mathscr C_0$-\Term{perfect-complex}{完美}的；
\item[ii)] \Term{canonical}{典范}\Term{morphism}{态射}（2.5.3）
\[
\operatorname{Proj}_0\circ\mathbf RF_k(L^\bullet)\longrightarrow L^\bullet
\]
是 $D^+(\mathscr A)$ 中的\Term{isomorphism}{同构}。
\end{enumerate}
\end{SGAresult}
\begin{proof}
这里所考虑的 $\mathscr C_0$-\Term{perfect-complex}{完美复形}概念是\Term{local-scheme}{局部}的，即相对于 $\mathscr C_0$ 所定义的 $S$ 上的\Term{fibered-category}{纤维范畴}而言的\Term{perfect}{完美}性。这个\Term{perfect}{完美}性概念确实有意义，因为 $\mathscr C_0$ 满足第 I 讲 4.0 的公理；其中只有对\Term{epimorphism}{满态射}的\Term{kernel}{核}的\Term{stable}{稳定}性不是直接可见的，而这由 1.5.1 得到。

因 $S$ \Term{quasi-compact}{拟紧}，只须在 $S$ 上\Term{local-scheme}{局部}考察 2.6 的两个断言，故可设 $S$ \Term{affine}{仿射}。此时\Term{invertible-sheaf}{可逆层} $\mathcal O_X(1)$ 为\Term{ample}{丰沛}，且 $X$ 上存在全局\Term{resolution}{消解}（第 II 讲 2.2.9）。因此，存在 $D_{\mathrm{parf}}(\mathscr A)$ 中的\Term{strictly-perfect}{严格完美}\Term{complex}{复形} $L'^\bullet$，以及 $D_{\mathrm{parf}}(\mathscr A)$ 中的\Term{isomorphism}{同构} $L^\bullet\xleftarrow{\sim}L'^\bullet$，从而可设 $L^\bullet$ \Term{strictly-perfect}{严格完美}。但在 2.4 中已经看到，当 $k$ 充分大时，除\Term{finite}{有限}多个以外均为零的各 $L^i$ 都是 $F_k$-\Term{acyclic-object}{无上同调}的。因此，$k$ 充分大时有
\[
\mathbf RF_k(L^\bullet)=F_k(L^\bullet).
\]
由 1.3 和 1.6，可取 $k_0$，使 $k\geq k_0$ 时既满足上面的条件，又使各 $F_k(L^i)$ 均为 $\mathscr C_0$-\Term{perfect-complex}{完美}的 $B$-\Term{module}{模}。于是\Term{complex}{复形} $F_k(L^\bullet)$ 是 $\mathscr C_0$-\Term{perfect-complex}{完美}的，i）得证。

若 $k$ 按上述方式选取，则\Term{homomorphism}{同态}（2.5.3）就是\Term{homomorphism}{同态}
\SourcePage{397}
\[
\operatorname{Proj}_0(F_k(L^\bullet))\longrightarrow L^\bullet.
\]
因已设 $L^\bullet$ \Term{strictly-perfect}{严格完美}，它在每个\Term{set-component}{分量}上都是\Term{isomorphism}{同构}，故 ii）成立。
\end{proof}
\textbf{2.7。} 在 $\mathscr C$ 上，将\Term{graded-module}{分次} $B$-\Term{module}{模}对应到它的第 $n$ 个\Term{set-component}{分量}的\Term{functor}{函子}，是取值于 $\mathscr C'$ 的\Term{exact-functor}{正合函子}，因此定义\Term{functor}{函子}
\[
\operatorname{Comp}_n:D^+(\mathscr C)\longrightarrow D^+(\mathscr C').\tag{2.7.1}
\]
对每个 $n$，可考虑由与 $B(-n)$ 作\Term{tensor-product}{张量积}定义的\Term{functor}{函子} $\mathscr C'\to\mathscr C$。\Term{homomorphism}{同态}（1.2.1）在 $\mathscr C$ 上定义一个\Term{functor-morphism}{函子态射}，从而给出\Term{derived-category}{导出范畴}间相应\Term{functor}{函子}的一个\Term{morphism}{态射}：
\[
\operatorname{Comp}_n\otimes_{\mathcal O_S}B(-n)
\longrightarrow\operatorname{Id}_{D^+(\mathscr C)}.\tag{2.7.2}
\]
设 $M^\bullet$ 为 $\mathscr C_0$-\Term{perfect-complex}{完美复形}。下面证明存在一串\Term{distinguished-triangle}{区别三角}
\[
\begin{tikzcd}[column sep=small]
&P_1^\bullet\arrow[dl]&\\
N_1^\bullet\arrow[rr]&&M^\bullet\arrow[ul]
\end{tikzcd}
\qquad
\begin{tikzcd}[column sep=small]
&P_2^\bullet\arrow[dl]&\\
N_2^\bullet\arrow[rr]&&P_1^\bullet\arrow[ul]
\end{tikzcd}
\qquad\cdots\qquad
\begin{tikzcd}[column sep=small]
&P_k^\bullet\arrow[dl]&\\
N_k^\bullet\arrow[rr]&&P_{k-1}^\bullet\arrow[ul]
\end{tikzcd}
\]
（中间依次类推），其中 $P_k^\bullet$ 为\Term{acyclic-complex}{无上同调复形}，各 $N_i^\bullet$ 形如 $K_i^\bullet\otimes_{\mathcal O_S}B(d)$，其中 $K_i^\bullet$ 是 $S$ 上通常意义的\Term{perfect-complex}{完美复形}。

因 $M^\bullet$ 是 $\mathscr C_0$-\Term{perfect-complex}{完美}的，且 $S$ \Term{quasi-compact}{拟紧}，存在 $S$ 的一个\Term{finite}{有限}\Term{covering}{覆盖}及整数 $d_0$，使在此\Term{covering}{覆盖}的每个\Term{open-subset}{开集} $U$ 上，$M^\bullet$ \Term{isomorphism}{同构}于一个\Term{bounded-complex}{有界复形} $M_U^\bullet$，其各项是\Term{finite-type}{有限型}\Term{locally-free}{局部自由} $B$-\Term{module}{模}，并且 $d<d_0$ 时 $\operatorname{Comp}_d(M_U^\bullet)=0$。
\SourcePage{398}
现考虑\Term{canonical}{典范}\Term{morphism}{态射}（2.7.2）：
\[
N_1^\bullet=\operatorname{Comp}_{d_0}(M^\bullet)\otimes_{\mathcal O_S}B(-d_0)
\longrightarrow M^\bullet.
\]
令 $P_1^\bullet$ 为此\Term{morphism}{态射}的\Term{cone}{锥}。它仍然是 $\mathscr C_0$-\Term{perfect-complex}{完美复形}。事实上，对每个 $n$，$\operatorname{Comp}_n(M^\bullet)$ 是 $S$ 上的\Term{perfect-complex}{完美复形}，因为当 $M^\bullet$ 为各项是\Term{finite-type}{有限型}\Term{locally-free}{局部自由} $B$-\Term{module}{模}的\Term{bounded-complex}{有界复形}时确实如此；所以\Term{complex}{复形} $\operatorname{Comp}_n(M^\bullet)\otimes_{\mathcal O_S}B(-d_0)$ 是 $\mathscr C_0$-\Term{perfect-complex}{完美}的，其\Term{cone}{锥}也如此（第 I 讲 4.10）。

此外，在此\Term{covering}{覆盖}的每个\Term{open-subset}{开集} $U$ 上可写出\Term{commutative-diagram}{交换图}
\[
\begin{tikzcd}
\operatorname{Comp}_{d_0}(M^\bullet)\otimes_{\mathcal O_S}B(-d_0)
\arrow[r]\arrow[d,"\sim"']&M^\bullet\\
\operatorname{Comp}_{d_0}(M_U^\bullet)\otimes_{\mathcal O_S}B(-d_0)
\arrow[r]&M_U^\bullet\arrow[u,"\sim"'],
\end{tikzcd}
\]
它在两\Term{cone}{锥}之间诱导\Term{isomorphism}{同构}。但下面的\Term{homomorphism}{同态}在\Term{degree}{次数} $d_0$ 是\Term{isomorphism}{同构}，故 $P_1^\bullet$ 在 $U$ 上\Term{isomorphism}{同构}于一个各项为\Term{finite-type}{有限型}\Term{locally-free}{局部自由} $B$-\Term{module}{模}的\Term{bounded-complex}{有界复形}，且其\Term{set-component}{分量} $\operatorname{Comp}_d$ 在 $d<d_0+1$ 时均为零。因此，$P_1^\bullet$ 满足 $M^\bullet$ 的同样性质，只是 $d_0$ 被 $d_0+1$ 替代。于是可迭代这一过程；只须说明，\Term{finite}{有限}次操作后所得\Term{complex}{复形} $P_k^\bullet$ 为\Term{acyclic-complex}{无上同调复形}。

这仍是一个\Term{local-scheme}{局部}问题，故只须对 $M_U^\bullet$ 考察。而对每一项 $M_U^i$，上述各\Term{morphism}{态射}正是定义命题 1.7 中\Term{filtration}{滤过}的那些\Term{morphism}{态射}；此\Term{filtration}{滤过}是\Term{finite}{有限}的，所以\Term{finite}{有限}次操作后，就从 $M_U^\bullet$ 得到零\Term{complex}{复形}，即相应的\Term{complex}{复形} $P_k^\bullet$ 是\Term{acyclic-complex}{无上同调}的。
\SourcePage{399}
\textbf{2.8。} 设 $L^\bullet$ 是 $X$ 上的\Term{perfect-complex}{完美复形}。存在整数 $k$，使 $M^\bullet=\mathbf RF_k(L^\bullet)$ 是 $\mathscr C_0$-\Term{perfect-complex}{完美复形}，且 $L^\bullet$ \Term{isomorphism}{同构}于 $\operatorname{Proj}_0(\mathbf RF_k(L^\bullet))$。将 2.7 应用于\Term{complex}{复形} $M^\bullet$，再应用\Term{functor}{函子} $\operatorname{Proj}_0$，所构造的\Term{distinguished-triangle}{区别三角}序列便给出 $D_{\mathrm{parf}}(\mathscr A)$ 中的一串\Term{distinguished-triangle}{区别三角}（显然，$\operatorname{Proj}_0$ 将 $\mathscr C_0$-\Term{perfect-complex}{完美复形}变成\Term{perfect-complex}{完美复形}）。此外，对 $D^+(\mathscr C')$ 中的 $K^\bullet$ 有
\[
\operatorname{Proj}_0(K^\bullet\otimes_{\mathcal O_S}B(d))=f^*(K^\bullet)(d)
\]
（注意，由 $f$ 的\Term{flatness}{平坦性}，$f^*=\mathbf Lf^*$）。

于是，在 $K^\bullet(X)$ 中得到等式
\[
\operatorname{cl}(L^\bullet)
=\sum_i\operatorname{cl}(f^*(K_i^\bullet))\cdot\xi^{d_i}
=\sum_i f^*(\operatorname{cl}(K_i^\bullet))\cdot\xi^{d_i},
\]
其中各 $K_i^\bullet$ 是 2.7 中定义的 $S$ 上的\Term{perfect-complex}{完美复形}。因此，$\xi$ 的幂作为 $K^\bullet(S)$-\Term{module}{模}的\Term{generator}{生成元}生成 $K^\bullet(X)$。

要证明 $1,\xi,\ldots,\xi^{r+1}$ \Term{linearly-dependent}{线性相关}，考虑\Term{homomorphism}{同态}
\[
K^\bullet_{\text{naïf}}(X)\longrightarrow K^\bullet_{\mathrm{parf}}(X),
\]
它将\Term{finite-type}{有限型}\Term{locally-free-sheaf}{局部自由层} $L$ 的类映到\Term{degree}{次数} $0$ 处等于 $L$、其他\Term{degree}{次数}处为零的\Term{complex}{复形}的类。于是方程（1.9.1）给出 $K^\bullet_{\mathrm{parf}}(X)$ 中的一个相关关系。

要证明 $1,\xi,\ldots,\xi^r$ \Term{linearly-independent}{线性无关}，可使用 $K^\bullet(S)$-\Term{linear}{线性}\Term{morphism}{态射} $f_*:K^\bullet(X)\to K^\bullet(S)$（第 IV 讲 2.11.1）。有
\[
f_*(1)=1,
\]
因为 $f_*(\mathcal O_X)=\mathcal O_S$，且 $i\geq1$ 时 $R^i(f_*(\mathcal O_X))=0$；又有
\[
f_*(\xi^n)=0\qquad(-r\leq n<0),
\]
因为 $i\geq0$ 时 $R^i(f_*(\mathcal O_X(n)))=0$。

若有形如
\[
f^*(x_0)+f^*(x_1)\xi+\cdots+f^*(x_r)\xi^r=0
\]
的关系，将其乘以 $\xi^{-r}$ 并应用 $f_*$，便得 $x_r=0$。
\SourcePage{400}
再用归纳完成证明：注意，若 $0\leq s\leq r$，且 $s<j\leq r$ 时 $x_j=0$，则有
\[
\sum_{i=0}^s f^*(x_i)\xi^i=0,
\]
从而
\[
\sum_{i=0}^s f^*(x_i)\xi^{1+r-s}=0,
\]
由上面的论述即得 $x_s=0$。
